\documentclass[11pt]{article}

% ---------- Packages ----------
\usepackage[margin=1in]{geometry}
\usepackage{amsmath, amssymb, amsthm}
\usepackage{mathtools}
\usepackage{enumitem}
\usepackage{fancyhdr}
\usepackage{xcolor}
\usepackage{hyperref}
\hypersetup{colorlinks=true, linkcolor=blue, citecolor=blue, urlcolor=blue}

% ---------- Header / Footer ----------
\setlength{\headheight}{14pt}
\pagestyle{fancy}
\fancyhf{}
\lhead{MA521 --- Symmetry, Geometry, and Representations}
\rhead{Homework 1}
\cfoot{\thepage}
\renewcommand{\headrulewidth}{0.4pt}

% ---------- Theorem-like environments ----------
\theoremstyle{plain}
\newtheorem{theorem}{Theorem}[section]
\newtheorem{lemma}[theorem]{Lemma}
\newtheorem{corollary}[theorem]{Corollary}
\newtheorem{proposition}[theorem]{Proposition}

\theoremstyle{definition}
\newtheorem{definition}[theorem]{Definition}

% ---------- Problem / Solution environment ----------
% \begin{problem}[optional label] ... \end{problem} typesets the given
% statement; \begin{solution} ... \end{solution} typesets the write-up
% and appends a QED-style marker at the end.
\newcounter{problemnum}
\newenvironment{problem}[1][]{%
  \refstepcounter{problemnum}%
  \par\medskip\noindent\textbf{Problem \theproblemnum{#1}.}\ \itshape
}{\par\medskip}

\newenvironment{solution}{%
  \par\noindent\textit{Solution.}\
}{%
  \hfill$\blacksquare$\par\bigskip
}

% ---------- Common shorthands ----------
\newcommand{\Hom}{\operatorname{Hom}}
\newcommand{\Aut}{\operatorname{Aut}}
\newcommand{\End}{\operatorname{End}}
\newcommand{\Iso}{\operatorname{Iso}}
\newcommand{\Span}{\operatorname{span}}
\newcommand{\Dim}{\operatorname{dim}}
\newcommand{\R}{\mathbb{R}}
\newcommand{\Q}{\mathbb{Q}}
\newcommand{\Z}{\mathbb{Z}}
\newcommand{\C}{\mathbb{C}}
\newcommand{\F}{\mathbb{F}}
\newcommand{\dual}[1]{#1^{\ast}}
\newcommand{\Herm}{\mathcal{H}}

\DeclareMathOperator{\SU}{SU}
\DeclareMathOperator{\SO}{SU}
\DeclareMathOperator{\GL}{GL}
\DeclareMathOperator{\SL}{SL}
\DeclareMathOperator{\pf}{pf}

\title{\vspace{-2em}MA521: Symmetry, Geometry, and Representations
\\ \large Homework 3}
\author{Asher Lucas-Cuddeback}

\begin{document}
\maketitle
\thispagestyle{fancy}

% Start numbering at 6 so the first problem below is Problem 6.
\setcounter{problemnum}{5}

\begin{problem}
  Consider the following function $H \colon V \times V \to \C$ on
  $V = \C^2$:
  \[
    H\!\left(
      \begin{bmatrix} z_1 \\ z_2
      \end{bmatrix},
      \begin{bmatrix} z_3 \\ z_4
    \end{bmatrix}\right)
    = 2\overline{z_1} z_3 - i\,\overline{z_1} z_4
    + i\,\overline{z_2} z_3 + 2\,\overline{z_2} z_4.
  \]
  \begin{enumerate}[label=(\alph*)]
    \item Explain why $H$ is a Hermitian form, and write $H$ in the form
      $H(v,w) = v^{*} B w$, where $B$ is some $2 \times 2$ Hermitian
      matrix.

    \item Show that the underlying signature of this Hermitian form is
      $(2,0)$, i.e.\ it is equivalent to the standard Hermitian form on
      $\C^2$ given by:
      \[
        \Herm\!\left(
          \begin{bmatrix} z_1 \\ z_2
          \end{bmatrix},
          \begin{bmatrix} z_3 \\ z_4
        \end{bmatrix}\right)
        = \overline{z_1} z_3 + \overline{z_2} z_4.
      \]

    \item Try to write down a Hermitian form on $\C^2$ that is not equal
      to, but equivalent to, the signature $(1,1)$ form given by:
      \[
        \mathcal{H}\!\left(
          \begin{bmatrix} z_1 \\ z_2
          \end{bmatrix},
          \begin{bmatrix} z_3 \\ z_4
        \end{bmatrix}\right)
        = \overline{z_1} z_3 - \overline{z_2} z_4.
      \]
  \end{enumerate}
\end{problem}

\begin{solution}
  \begin{proof}[Proof of 6(a)]
    To show that this form is Hermitian, it suffices to show that
    \[
      \Herm\left(
        \begin{bmatrix}
          z_1 \\
          z_2
        \end{bmatrix},
        \begin{bmatrix}
          z_3 \\
          z_4
      \end{bmatrix}\right)=\overline{
        \Herm\left(
          \begin{bmatrix}
            z_3 \\
            z_4
          \end{bmatrix},
          \begin{bmatrix}
            z_1 \\
            z_2
        \end{bmatrix}\right)
      }
    \]
    We are given:
    \[
      \Herm\!\left(
        \begin{bmatrix} z_1 \\ z_2
        \end{bmatrix},
        \begin{bmatrix} z_3 \\ z_4
      \end{bmatrix}\right)
      = 2\overline{z_1} z_3 - i\,\overline{z_1} z_4
      + i\,\overline{z_2} z_3 + 2\,\overline{z_2} z_4.
    \]
    Thus
    \[
      \Herm\left(
        \begin{bmatrix}
          z_3 \\
          z_4
        \end{bmatrix},
        \begin{bmatrix}
          z_1 \\
          z_2
      \end{bmatrix}\right)= 2\overline{z_3}z_1 - i\overline{z_3}z_2
      + i\overline{z_4}z_1 + 2\overline{z_4}z_2
    \]
    Taking the Hermetian conjugate of this we see that
    \[
      \overline{2\overline{z_3}z_1 - i\overline{z_3}z_2
      + i\overline{z_4}z_1 + 2\overline{z_4}z_2} =
      2\overline{z_1}z_3 + i\overline{z_2}z_3 - i\overline{z_1}z_4 +
      2\overline{z_2}z_4
    \]
    Which we can rearrange to obtain our desired result
    \[
      = 2\overline{z_1} z_3 - i\,\overline{z_1} z_4
      + i\,\overline{z_2} z_3 + 2\,\overline{z_2} z_4.
    \]
    The matrix describing this Hermetian form is
    \[
      \Herm\left(
        \begin{bmatrix}
          z_1 \\
          z_2
        \end{bmatrix},
        \begin{bmatrix}
          z_3 \\
          z_4
      \end{bmatrix}\right)=
      \begin{bmatrix}
        z_1 & z_2
      \end{bmatrix}^*
      \begin{bmatrix}
        2 & -i \\
        i & 2
      \end{bmatrix}
      \begin{bmatrix}
        z_3 \\
        z_4
      \end{bmatrix}
    \]
  \end{proof}

  \begin{proof}[Proof of 6(b)]
    We can achieve this with a mildly modified version of the
    'perform row op and the same column op' method for Bilinear
    forms. The modification is that we will do the Hermetian
    conjugate row op and then the column op. So
    \[
      \begin{bmatrix}
        2 & -i \\
        i & 2
      \end{bmatrix}\xrightarrow[iC_2]{-iR_2}
      \begin{bmatrix}
        2 & 1 \\
        1 & 2
      \end{bmatrix}\xrightarrow[2R_2 - R_1]{2C_2-C_1}
      \begin{bmatrix}
        2 & 0 \\
        0 & 5
      \end{bmatrix}\xrightarrow[\frac{1}{\sqrt{2}}R_1]{\frac{1}{\sqrt{2}}C_1}
      \begin{bmatrix}
        1 & 0 \\
        0 & 5
      \end{bmatrix}\xrightarrow[\frac{1}{\sqrt{5}}R_2]{\frac{1}{\sqrt{5}}C_2}
      \begin{bmatrix}
        1 & 0 \\
        0 & 1
      \end{bmatrix}
    \]
    Which is the matrix describing
    \[
      \Herm\!\left(
        \begin{bmatrix} z_1 \\ z_2
        \end{bmatrix},
        \begin{bmatrix} z_3 \\ z_4
      \end{bmatrix}\right)
      = \overline{z_1} z_3 + \overline{z_2} z_4.
    \]
  \end{proof}

  \begin{proof}[Proof of 6(c)]
    The matrix
    \[
      A=
      \begin{bmatrix}
        1 & 0 \\
        0 & -1
      \end{bmatrix}
    \]
    is the Hermetian form for which any other signature \( (1, 1) \)
    Hermitian form must be congruent to. So we can simply do the same
    row-column operation trick with the reverse intent. So
    \[
      \begin{bmatrix}
        1 & 0 \\
        0 & -1
      \end{bmatrix}\xrightarrow[R_2 + R_1]{C_2 + C_1}
      \begin{bmatrix}
        1 & 1 \\
        1 & -1
      \end{bmatrix}
    \]
    satisfies the requirements (if in an unsatisfyingly low-effort
    manner I'm sorry).
  \end{proof}
\end{solution}

\begin{problem}
  Give explicit equational descriptions of the following matrix groups.
  \begin{enumerate}[label=(\alph*)]
    \item $\SU(2)$
    \item $\SU(1,1)$
    \item $\SO_2(\C)$
  \end{enumerate}
\end{problem}

\begin{solution}
  \begin{proof}[Proof of 7(a)]
    The group \( \SU(2) \) is the Unitary matrices with determinant 1. So
    \[
      \SU(2)=\left\{M\in \GL_2(\C)\,|\,\det(M)=1 \text{ and } M^*M=I\right\}
    \]
    \[
      =\left\{
        \begin{bmatrix}a & b \\ c & d
        \end{bmatrix}\in \SL_2(\C)\,|\,
        \begin{bmatrix}
          \overline{a} &
          \overline{c} \\
          \overline{b} &
          \overline{d}
        \end{bmatrix}
        \begin{bmatrix}
          a & b \\ c & d
        \end{bmatrix}=
        \begin{bmatrix}
          1 & 0 \\
          0 & 1
        \end{bmatrix}
      \right\}
    \]
    \[
      =\left\{
        \begin{bmatrix}a & b \\ c & d
        \end{bmatrix}\in \SL_2(\C)\,|\,
        \begin{bmatrix}
          \overline{a}a + \overline{c}c & \overline{a}b +
          \overline{c}d \\ \overline{b}a + \overline{d}c &
          \overline{b}b + \overline{d}d
        \end{bmatrix}=
        \begin{bmatrix}
          1 & 0 \\
          0 & 1
        \end{bmatrix}
      \right\}
    \]
    This can be thought of as the matrix derived from the vectors
    \(
      \begin{bmatrix}
        a \\ b
    \end{bmatrix} \) and \(
      \begin{bmatrix}
        c \\ d
    \end{bmatrix} \)
    And
    \[
      \Herm\!\left(v_1,v_2\right) = \overline{a} c + \overline{b} d
    \]
    Where we have the matrix
    \[
      \begin{bmatrix}
        \Herm(v_1, v_1) & \Herm(v_1, v_2) \\
        \Herm(v_2, v_1) & \Herm(v_2, v_2)
      \end{bmatrix}.
    \]
    This representation implies that \( v_1 \) and \( v_2 \) are
    linearly independent (since \( \Herm(v_1,v_2)=0 \))
  \end{proof}

  \begin{proof}[Proof of 7(b)]
  \end{proof}

  \begin{proof}[Proof of 7(c)]
  \end{proof}
\end{solution}

\begin{problem}
  \begin{enumerate}[label=(\alph*)]
    \item Learn a bit about the Pfaffian of a skew-symmetric matrix so
      that you can demonstrate that any matrix in $Sp_{2n}(F)$
      necessarily has determinant $1$.

    \item Argue that $Sp_{2n}(F)$ is always a subgroup of $SL_{2n}(F)$.

    \item In class we saw that $Sp_2(F) = SL_2(F)$. Show that
      $SL_{2n}(F)$ is strictly larger than $Sp_{2n}(F)$ for $n \geq 2$.
  \end{enumerate}
\end{problem}

\begin{solution}
  \begin{proof}[Proof of 8(a)]
    The Pfaffian has the property \( \pf(A)^2 = \det A \)
  \end{proof}

  \begin{proof}[Proof of 8(b)]
    Since \( \forall A \in Sp_{2n}(F)\), \( \det A = 1 \) we are
    naturally a subset of \( \SL_{2n}(F) \). Additionally, since \(
    \det A \neq 0 \), \( A^{-1}\in Sp_{2n}(F) \) and finally, since
    the determinant is multiplicative, \( \forall A,B\in
    Sp_{2n}(F),\,\det (AB)=1 \) so of course \( AB\in Sp_{2n}(F) \).
  \end{proof}

  \begin{proof}[Proof of 8(c)]
  \end{proof}
\end{solution}

\end{document}
